\clearpage
\chapter*{Erratum à l'exposé n\textsuperscript{o} 182}
\addcontentsline{toc}{chapter}{Erratum à l'exposé n\textsuperscript{o} 182}
\setcounter{footnote}{1}
\src{298}

\noindent\textbf{2. Géométrie formelle et géométrie algébrique.}

\begin{center}
[Séminaire Bourbaki, t. 11, 1958/59, n\textsuperscript{o} 182, 28 p.]
\end{center}

\noindent\textbf{Page 8, théorème 6.} --- L'hypothèse « $f$
quasi-projectif » peut être remplacée par l'hypothèse plus faible « $f$
séparé », en vertu du résultat suivant (cf. SGA VIII, 6.2) : tout morphisme
$f:X\longrightarrow Y$, qui est quasi-fini et séparé, est quasi-projectif.

\src{299}
\noindent\textbf{Page 12, et page 14, remarque 1.} --- Signalons que J.-P.
Serre a construit une variété projective non singulière, de dimension 3, sur
un corps algébriquement clos $k$, de caractéristique $p>0$, qui ne provient pas
par réduction d'un schéma propre sur un anneau local intègre de corps résiduel
$k$, et ayant un corps de fractions de caractéristique
nulle.\footnote{Serre (Jean-Pierre). --- Exemples de variétés projectives en
caractéristique $p$ non relevables en caractéristique zéro, \emph{Proc. Nat.
Acad. Sc. U.S.A.}, t. 47, 1961, p. 108--109.} Mumford aurait trouvé un exemple
analogue, avec une surface projective non singulière.

\noindent\textbf{Page 15, remarques 2 et 3.} --- Je suis actuellement moins
optimiste concernant les résultats conjecturés ici. Cependant, le problème
relatif aux variations de structure pour l'espace projectif, signalé à la fin
de la remarque 3, a été résolu par l'affirmative par Hironaka, et le problème
analogue pour les variétés abéliennes a été résolu par Koizumi.

\noindent\textbf{Page 16, ligne 3.} --- Signalons que Mumford a construit
récemment les schémas des modules pour les courbes de genre $g$ (cf. SHMI). Le
théorème 10 prouve d'ailleurs que les « schémas de Jacobi d'échelon $n$ » de la
théorie des Modules sont non singuliers (et même simples sur $\mathbf Z$).

\noindent\textbf{Page 23, ligne 12, 1\textsuperscript{re} ligne du corollaire
5.} --- Ajouter : « $X$ ou $Y$ étant propre sur $k$ ».

La substance des §§ 1 à 5 est contenue dans la partie publiée de EGA III,
celle des §§ 6 et 7 est contenue dans SGA III. Pour l'étude du groupe
fondamental, voir SGA V, IX, X, XI, ainsi que SGA 1962 (exposés X, XII et XIII)
pour les théorèmes du type Lefschetz et de nombreuses questions ouvertes.
Seule la théorie des revêtements modérément ramifiés (cf. théorème 14) n'a pas
encore fait l'objet d'une rédaction en forme. Le corollaire du théorème 14, qui
donne la détermination complète des revêtements galoisiens d'ordre premier à la
caractéristique d'une courbe algébrique sur un corps algébriquement clos, a été
utilisé de façon essentielle à trois reprises :

1\textsuperscript{o} dans la démonstration par Igusa de l'inégalité de Picard
pour les surfaces projectives non singulières en caractéristique quelconque ;

2\textsuperscript{o} dans l'étude (développée indépendamment par Ogg et
Šafarevič) du groupe des fibrés principaux homogènes sous une variété abélienne
définie sur un corps de fonctions d'une variable, en caractéristique
quelconque ;

3\textsuperscript{o} dans la démonstration récente, par Artin, de certains
théorèmes-clefs sur la « cohomologie de Weil » des variétés algébriques.
